\chapter{Riproduzione dei risultati}
\label{app:riproducibilita}

\section{Build}

\begin{lstlisting}[language=shell,caption={Compilazione delle configurazioni usate nelle campagne.}]
make                                  # bin/matmul_mpi  (double, scheme_a)
make FORCE_GENERIC_K=1                # bin/matmul_mpi-generic
make PREC=float                       # bin/matmul_mpi-float
make KERNEL=cuda_naive                # richiede nvcc: solo sul server
make KERNEL=cuda_warp
make KERNEL=cuda_warp_smem
make KERNEL=cuda_warp_smem SMEM_PAD=0
make KERNEL=cuda_warp_tiled
make KERNEL=cublas                    # aggiunge -lcublas da solo
\end{lstlisting}

Ogni comando stampa in chiusura quale binario ha prodotto e con quali variabili.
Su una macchina priva di \code{nvcc} le build CUDA si fermano con un messaggio
esplicito e tutto il resto continua a funzionare.

\section{Validazione}

Da eseguire su taglie piccole, prima delle campagne e per ogni combinazione di
backend e precisione che si intende misurare.

\begin{lstlisting}[language=shell]
make test                     # funzioni indice, senza MPI
make check-cxx                # kernel.h e util.h includibili da C++
make check                    # prodotto distribuito contro oracolo, double
make PREC=float check         # stessa validazione in singola precisione
make check-padding            # lda = n_loc + 8
make KERNEL=cuda_warp check
make KERNEL=cuda_warp_smem check
make KERNEL=cublas check
\end{lstlisting}

\section{Esecuzione singola}

\begin{lstlisting}[language=shell]
# CPU, 4 rank su griglia 2x2
mpirun -np 4 --bind-to core --map-by core \
  ./bin/matmul_mpi -M 8000 -N 8000 -k 8 --pr 2 --pc 2 --reps 10

# GPU, rank singolo: la configurazione con cui si caratterizza un kernel
mpirun -np 1 --bind-to core --map-by core ./bin/matmul_mpi-cuda_warp \
  -M 20000 -N 20000 -k 8 --pr 1 --pc 1 --a-mode local --warmup 2 --reps 10
\end{lstlisting}

Il binding non \`e facoltativo sulle esecuzioni multi-rank: senza, le misure sono
rumore.

\section{Campagne}

\begin{lstlisting}[language=shell]
./script/run_cuda_experiments.sh --suite experiments/c3_cpu_baseline.txt
./script/run_cuda_experiments.sh --suite experiments/c2_k_sweep.txt
./script/run_cuda_experiments.sh --suite experiments/c1_size_ratio_cpu.txt
./script/run_cuda_experiments.sh --suite experiments/c1_size_ratio_gpu.txt
./script/run_cuda_experiments.sh --suite experiments/c6_grid_shape.txt
./script/run_cuda_experiments.sh --suite experiments/c4_strong_scaling.txt
./script/run_cuda_experiments.sh --suite experiments/c5_weak_scaling.txt
\end{lstlisting}

\begin{tcolorbox}[colback=todogiallo,colframe=todorosso,boxrule=0.5pt,
                  title=\textbf{Prima di rilanciare una campagna}]
Rilanciare sovrascrive i CSV della campagna. Archiviare la cartella:
\begin{lstlisting}[language=shell]
mv results/c4_strong_scaling results/c4_strong_scaling.$(date +%F_%H%M)
\end{lstlisting}
\end{tcolorbox}

\section{Unione dei CSV di una campagna}

Tutti i CSV aggregati di una cartella condividono l'intestazione:

\begin{lstlisting}[language=shell]
find results/c1_size_ratio_cpu -name '*.csv' \
     ! -name '*_raw.csv' ! -name '*_failures.csv' \
  | sort | xargs awk 'FNR==1 && NR!=1 {next} {print}' \
  > results/c1_size_ratio_cpu/merged.csv
\end{lstlisting}

I file \code{*\_raw.csv} hanno un'altra intestazione e vanno uniti separatamente.

\section{Note operative sul server}

\begin{itemize}[left=0pt]
  \item \textbf{Nodo di calcolo, non frontend.} Il nodo di login non ha accesso
        diretto alla GPU: \code{nvcc --version},
        \code{ompi\_info | grep cuda}, \code{deviceQuery} e \code{nvidia-smi}
        vanno eseguiti sui nodi di calcolo attraverso lo scheduler, e cos\`i tutti
        i benchmark.
  \item \textbf{Trasporto MPI.} Il runner usa \code{btl=self,sm} per default. Il
        componente \code{sm} non esiste in Open MPI 3.x/4.x, dove si chiama
        \code{vader}: se le esecuzioni multi-rank falliscono con un errore di
        BTL, verificare con \code{ompi\_info --param btl all} e aggiungere
        \code{--btl self,vader} sulla riga di comando, che ha la precedenza sul
        \code{.conf}.
  \item \textbf{bash $\ge$ 4.4.} Il runner espande array potenzialmente vuoti
        sotto \code{set -u}: su bash 3.2 (macOS) si ferma con
        \code{unbound variable}. Sul server Linux non si presenta.
  \item \textbf{Memoria della GPU.} Un blocco locale $\num{40000}^2$ in doppia
        precisione richiede \SI{12.8}{\gib} contro i \SI{15.5}{\gib} della
        scheda: il programma si ferma prima di allocare, indicando quanta memoria
        serve e quanta \`e libera.
\end{itemize}

\section{Compilazione di questa relazione}

\begin{lstlisting}[language=shell]
cd report && make          # oppure: latexmk -pdf main.tex
\end{lstlisting}

Per la versione di consegna, commentare la riga \code{\textbackslash todovisibiletrue}
in \file{preambolo.tex}: i marcatori dei dati mancanti spariscono dal PDF e
restano nel sorgente.
